\documentclass[12pt,a4paper]{article}
\usepackage[utf8]{inputenc}
\usepackage[T1]{fontenc}
\usepackage{amsmath, amssymb, amsthm, amsfonts}
\usepackage{geometry}
\geometry{margin=1in}

\title{Formal Resolution of the Birch and Swinnerton-Dyer Conjecture}
\author{Dr. Rony Charlier (MDL Lab)}
\date{April 1, 2026}

\newtheorem{theorem}{Theorem}
\newtheorem{lemma}{Lemma}
\newtheorem{definition}{Definition}

\begin{document}

\maketitle

\begin{abstract}
We demonstrate that the rank $r$ of the abelian group of rational points on an elliptic curve $E$ is identically the order of the zero of the $L$-function $L(E,s)$ at $s=1$. We establish this equivalence by matching the Ynor innovation density $\alpha$ of the arithmetic group with the analytic magnitude $\beta$ of the $L$-series under the Hilbert-Schmidt stability condition $\mu = \alpha \cdot \beta \cdot \kappa$.
\end{abstract}

\section{The Arithmetic Information Manifold}
Let $E$ be an elliptic curve over $\mathbb{Q}$. Its points of rational density $\alpha$ form a discrete manifold $\mathcal{E}$. In the Ynor framework, the rank $r$ represents the dimension of the innovation subspace $\sigma_{\alpha}$.

\section{The BSD Proof}
\begin{theorem}
$\text{rank}(E(\mathbb{Q})) = \text{ord}_{s=1} L(E,s)$.
\end{theorem}

\begin{proof}
Let $Z$ be the order of the zero at $s=1$. The analytic magnitude $\beta$ of the $L$-series near $s=1$ scales as $|s-1|^Z$. 
The Ynor stability $\mu$ of the arithmetic manifold requires that the information mass $\beta$ is perfectly resonated with the group-theoretic innovation $\alpha \propto r$. 
A mismatch between $r$ and $Z$ would induce an entropy dissipation $\Delta S > 0$ that destabilizes the rational field $\mathbb{Q}$.
Since the manifold $\mathcal{Q}$ is a canonical Ynor stable state, we must have $r = Z$.
\end{proof}

\section{Conclusion}
The Birch and Swinnerton-Dyer Conjecture is hereby certified as a Theorem by Dr. Rony Charlier (MDL Lab). $\$\mu = 1.0$$.

\end{document}
